\documentclass[14pt,a4paper]{extreport}

\usepackage{iftex}
\ifPDFTeX
  \usepackage[T2A]{fontenc}
  \usepackage[utf8]{inputenc}
  \usepackage{cmap}
  \usepackage{tempora}
\else
  \usepackage{fontspec}
  \setmainfont{Times New Roman}
  \setsansfont{Arial}
  \setmonofont{Courier New}
\fi
\usepackage[russian]{babel}
\usepackage[a4paper,left=30mm,top=20mm,right=10mm,bottom=20mm]{geometry}
\usepackage{setspace}
\onehalfspacing
\usepackage{indentfirst}
\setlength{\parindent}{1.25cm}
\setlength{\parskip}{0pt}
\usepackage{ragged2e}
\justifying
\sloppy
\emergencystretch=3em
\usepackage{caption}
\usepackage{array}
\usepackage{longtable}
\usepackage{tabularx}
\usepackage{booktabs}
\usepackage{enumitem}
\usepackage{titlesec}
\usepackage{tocloft}
\usepackage{chngcntr}
\usepackage{xurl}
\usepackage{listings}
\usepackage{hyperref}
\hypersetup{hidelinks,unicode=true}

\counterwithin{figure}{chapter}
\counterwithin{table}{chapter}
\renewcommand{\thefigure}{\thechapter.\arabic{figure}}
\renewcommand{\thetable}{\thechapter.\arabic{table}}
\DeclareCaptionLabelSeparator{aitudot}{ --- }
\captionsetup{font=normalsize,labelfont=normalfont,labelsep=aitudot}
\captionsetup[table]{position=top,justification=raggedright,singlelinecheck=false,name=Таблица,margin={1.25cm,0cm}}
\captionsetup[figure]{position=bottom,justification=centering,singlelinecheck=false,name=Рисунок}

\titleformat{\chapter}[block]{\normalfont\centering\fontsize{14pt}{21pt}\selectfont}{\thechapter}{1em}{\MakeUppercase}
\titlespacing*{\chapter}{0pt}{0pt}{12pt}
\titleformat{\section}[block]{\normalfont\centering\fontsize{14pt}{21pt}\selectfont}{\thesection}{1em}{}
\titlespacing*{\section}{0pt}{12pt}{6pt}

\newcommand{\unnumberedchapter}[1]{%
  \clearpage
  \phantomsection
  \addcontentsline{toc}{chapter}{#1}
  \begin{center}\fontsize{14pt}{21pt}\selectfont #1\end{center}\vspace{6pt}
}
\newcommand{\appsection}[1]{%
  \clearpage
  \phantomsection
  \addcontentsline{toc}{section}{#1}
  \begin{flushright}#1\end{flushright}
  \vspace{6pt}
}
\newcommand{\filledfield}[2]{\underline{\makebox[#1][l]{\hspace{0.15cm}#2}}}
\newcommand{\blankfield}[1]{\underline{\makebox[#1][l]{}}}
\lstset{
  basicstyle=\ttfamily\small,
  breaklines=true,
  columns=fullflexible,
  frame=single,
  xleftmargin=1.25cm,
  keepspaces=true
}

\renewcommand{\contentsname}{СОДЕРЖАНИЕ}
\addto\captionsrussian{\renewcommand{\contentsname}{СОДЕРЖАНИЕ}}
\renewcommand{\cfttoctitlefont}{\hfill\fontsize{14pt}{21pt}\selectfont}
\renewcommand{\cftaftertoctitle}{\hfill\vspace{6pt}}
\renewcommand{\cftchapfont}{\normalfont}
\renewcommand{\cftchappagefont}{\normalfont}
\renewcommand{\cftsecfont}{\normalfont}
\renewcommand{\cftsecpagefont}{\normalfont}
\renewcommand{\cftchapdotsep}{\cftdotsep}
\setlength{\cftbeforechapskip}{2pt}
\setlength{\cftbeforesecskip}{1pt}
\setcounter{tocdepth}{2}
\pagestyle{plain}

\begin{document}
\pagenumbering{arabic}

\begin{titlepage}
\thispagestyle{empty}
\begin{center}
Министерство просвещения Республики Казахстан\\[0.5cm]
Astana IT University College\\[0.7cm]
Специальность: 06130100 --- Программное обеспечение (по видам)

\vfill
{\bfseries ДИПЛОМНЫЙ ПРОЕКТ}\\[0.8cm]
Тема: Разработка веб-приложения ChessView
\vfill
\begin{flushright}
Студент: \filledfield{6cm}{Сакенов Ануар Сункарулы}\\[0.4cm]
Группа: \filledfield{6cm}{ПО 2302}\\[0.4cm]
Руководитель: \filledfield{6cm}{Мартынцов Николай Викторович}
\end{flushright}
\vfill
Астана, 2026
\end{center}
\end{titlepage}

\begin{titlepage}
\thispagestyle{empty}
\begin{center}
Министерство просвещения Республики Казахстан\\
Astana IT University College\\[1cm]
\end{center}
\begin{flushright}
Допущен к защите\\
Председатель ПЦК: \blankfield{7cm}\\
Дата: \blankfield{7cm}
\end{flushright}
\vfill
\begin{center}
{\bfseries ДИПЛОМНЫЙ ПРОЕКТ}\\[0.5cm]
на тему: «Разработка веб-приложения ChessView»\\[0.8cm]
по специальности 06130100 --- Программное обеспечение (по видам)
\end{center}
\vfill
\begin{flushright}
Выполнил: \filledfield{6cm}{Сакенов Ануар Сункарулы}\\[0.4cm]
Группа: \filledfield{6cm}{ПО 2302}\\[0.4cm]
Руководитель: \filledfield{6cm}{Мартынцов Николай Викторович}
\end{flushright}
\vfill
\begin{center}Астана, 2026\end{center}
\end{titlepage}

\begin{titlepage}
\thispagestyle{empty}
\begin{center}
{\bfseries ЗАДАНИЕ}\\[0.5cm]
на дипломный проект
\end{center}
Студент: \filledfield{8cm}{Сакенов Ануар Сункарулы}

Группа: \filledfield{8cm}{ПО 2302}

Тема дипломного проекта: «Разработка веб-приложения ChessView».

Номер приказа: \blankfield{7cm}

Дата выдачи задания: \blankfield{7cm}

\vspace{0.5cm}
Задание: разработать веб-приложение ChessView для шахматного онлайн-взаимодействия пользователей, включающее авторизацию, профиль, live-партию, серверную проверку ходов, WebSocket-события, историю партий, анализ позиций, задачи, турниры и документацию по запуску.

\vfill
\begin{flushright}
Руководитель: \blankfield{6cm}\\[0.8cm]
Студент: \blankfield{6cm}
\end{flushright}
\end{titlepage}

\clearpage
\setcounter{page}{4}
\tableofcontents

\unnumberedchapter{ВВЕДЕНИЕ}

Актуальность дипломного проекта ChessView связана с тем, что шахматные онлайн-сервисы объединяют несколько важных для разработки программного обеспечения задач: регистрацию пользователей, хранение профилей, проверку правил предметной области, работу с базой данных, обмен событиями в реальном времени и удобный интерфейс для анализа партий. Такая предметная область позволяет показать не только отдельную страницу с шахматной доской, но и полноценное full-stack приложение учебного уровня.

Практическая значимость проекта состоит в создании расширяемой шахматной платформы, пригодной для демонстрации ключевых механизмов современной веб-разработки. В проекте реализованы клиентская часть на React и TypeScript, backend на FastAPI, хранение данных в PostgreSQL, миграции Alembic, REST API, WebSocket endpoint, серверная проверка ходов через python-chess, локальный анализ позиций с использованием Stockfish worker, задачи, турниры и административный интерфейс.

Целью дипломного проекта является разработка веб-приложения ChessView, предназначенного для организации шахматного онлайн-взаимодействия пользователей, ведения профилей, проведения live-партий, анализа позиций и демонстрации связанных учебных модулей платформы.

Для достижения цели были определены следующие задачи: изучить особенности шахматных веб-платформ; рассмотреть аналоги Chess.com, Lichess и Tornelo; определить функциональные и нефункциональные требования; выбрать стек технологий; спроектировать архитектуру приложения; реализовать backend, frontend и базу данных; реализовать WebSocket-сценарии live-игры; подготовить документацию запуска и проверки; описать ограничения текущей реализации.

Объектом разработки является шахматное веб-приложение. Предметом разработки являются программные модули ChessView, обеспечивающие авторизацию, игровой процесс, хранение данных, real-time события, анализ позиций, задачи, турниры и расширенные пользовательские сценарии. Результатом работы является репозиторий ChessView по адресу \url{https://github.com/thefueki/chessview}, документация и пояснительная записка.

\chapter{ОСНОВНАЯ ЧАСТЬ}

\section{Общие понятия веб-приложений и шахматных онлайн-платформ}

Веб-приложение представляет собой программную систему, с которой пользователь работает через браузер. В отличие от статического сайта, веб-приложение содержит состояние, пользовательские действия, сетевой обмен и бизнес-логику. Для ChessView это означает наличие frontend-интерфейса, backend API, WebSocket-канала и базы данных.

Шахматная онлайн-платформа предъявляет дополнительные требования. Она должна корректно отображать доску, соблюдать правила игры, хранить историю ходов, учитывать контроль времени и синхронизировать состояние между двумя игроками. В ChessView live-партия построена по server-authoritative модели: клиент отправляет ход, но окончательное изменение состояния выполняется на сервере.

REST API используется для операций типа запрос-ответ: регистрация, вход, загрузка профиля, история партий, задачи и турниры. WebSocket используется для событий, где важна быстрая доставка: матчмейкинг, ход, состояние партии, чат и WebRTC signaling. Такое разделение упрощает архитектуру и соответствует фактической структуре проекта.

\section{Обзор существующих решений}

При анализе аналогов важно учитывать масштаб. Chess.com и Lichess являются зрелыми массовыми платформами, а Tornelo ориентирован на организацию турниров. ChessView не позиционируется как конкурент этим сервисам. Его назначение --- учебная full-stack реализация ключевых механизмов шахматной платформы.

\begin{center}
\scriptsize
\setlength{\tabcolsep}{3pt}
\renewcommand{\arraystretch}{1.25}
\begin{longtable}{|>{\RaggedRight\arraybackslash}p{0.14\linewidth}|>{\RaggedRight\arraybackslash}p{0.20\linewidth}|>{\RaggedRight\arraybackslash}p{0.24\linewidth}|>{\RaggedRight\arraybackslash}p{0.26\linewidth}|}
\caption{Сравнение аналогов и ChessView}\label{tab:analogues}\\
\hline
Решение & Назначение & Характерные функции & Что учтено в ChessView\\ \hline
\endfirsthead
\hline
Решение & Назначение & Характерные функции & Что учтено в ChessView\\ \hline
\endhead
Chess.com & Массовая шахматная платформа & Игра, рейтинги, задачи, обучение, турниры, профиль & Связь игры, профиля, истории, задач и рейтингов\\ \hline
Lichess & Бесплатная шахматная платформа & Быстрый доступ к игре, анализ, задачи, турниры & Разделение live-игры и анализа, локальная работа с позициями\\ \hline
Tornelo & Турнирная платформа & Регистрация, паринги, расписания, OTB-сценарии & Турнирный модуль, Swiss helpers, OTB tournament type\\ \hline
ChessView & Учебное full-stack приложение & Auth, live game, WebSocket, анализ, задачи, турниры, admin & Демонстрация архитектуры и ключевых механизмов без заявления о промышленном масштабе\\ \hline
\end{longtable}
\end{center}

\section{Требования к ChessView}

Функциональные требования приведены в таблице~\ref{tab:functional}.

\begin{center}
\scriptsize
\setlength{\tabcolsep}{4pt}
\renewcommand{\arraystretch}{1.25}
\begin{longtable}{|>{\RaggedRight\arraybackslash}p{0.08\linewidth}|>{\RaggedRight\arraybackslash}p{0.28\linewidth}|>{\RaggedRight\arraybackslash}p{0.46\linewidth}|>{\RaggedRight\arraybackslash}p{0.10\linewidth}|}
\caption{Функциональные требования}\label{tab:functional}\\
\hline
Код & Требование & Описание & Статус\\ \hline
\endfirsthead
\hline
Код & Требование & Описание & Статус\\ \hline
\endhead
F1 & Авторизация & Регистрация, вход, refresh token, текущий пользователь, protected routes & Реализовано\\ \hline
F2 & Профиль & Рейтинг, статистика, аватар, история, leaderboard, сравнение игроков & Реализовано\\ \hline
F3 & Матчмейкинг & Очередь через WebSocket и создание партии с контролем времени & Реализовано, in-memory\\ \hline
F4 & Live-партия & Ходы, часы, результат, reconnect, timeout, resign/draw & Реализовано\\ \hline
F5 & Анализ & FEN/PGN, board editor, локальный Stockfish worker & Реализовано\\ \hline
F6 & Задачи & Каталог задач и учет попыток пользователя & Реализовано\\ \hline
F7 & Турниры & Создание, участие, standings, Swiss helpers, round progression & Реализовано\\ \hline
F8 & Расширенные модули & Shop, clubs, OTB, scheduled matches, admin, verification & Частично: часть модулей демонстрационная\\ \hline
\end{longtable}
\end{center}

Нефункциональные требования приведены в таблице~\ref{tab:nonfunctional}.

\begin{center}
\scriptsize
\setlength{\tabcolsep}{4pt}
\renewcommand{\arraystretch}{1.25}
\begin{longtable}{|>{\RaggedRight\arraybackslash}p{0.08\linewidth}|>{\RaggedRight\arraybackslash}p{0.30\linewidth}|>{\RaggedRight\arraybackslash}p{0.42\linewidth}|>{\RaggedRight\arraybackslash}p{0.12\linewidth}|}
\caption{Нефункциональные требования}\label{tab:nonfunctional}\\
\hline
Код & Требование & Описание & Проверка\\ \hline
\endfirsthead
\hline
Код & Требование & Описание & Проверка\\ \hline
\endhead
N1 & Разделение слоев & Отдельные frontend, backend, admin frontend и PostgreSQL & Структура репозитория\\ \hline
N2 & Серверная проверка ходов & Клиент не определяет итоговое состояние live-партии & `backend/domains/game/domain/moves.py`\\ \hline
N3 & Миграции & Схема базы управляется Alembic & `backend/alembic/versions`\\ \hline
N4 & Поддерживаемость & Домены backend и feature-sliced структура frontend & Структура директорий\\ \hline
N5 & Защита API & Protected endpoints требуют JWT & HTTP 401 без токена\\ \hline
N6 & Воспроизводимый запуск & Docker Compose и split dev workflow & README, `docker-compose.yml`\\ \hline
\end{longtable}
\end{center}

\section{Обоснование технологий}

React и TypeScript выбраны для компонентного интерфейса с типизированными API-ответами. Vite обеспечивает быстрый dev server и сборку. TanStack Query используется для server state, Zustand --- для локального состояния пользователя, матчмейкинга и UI.

FastAPI выбран для typed REST API, зависимостей и WebSocket support. SQLAlchemy async и PostgreSQL обеспечивают реляционное хранение пользователей, партий, ходов, задач, турниров и платежных записей. Alembic используется для версионирования схемы. Python-chess применяется на backend для проверки ходов, а chess.js и Stockfish worker используются на frontend для анализа и учебных сценариев.

\chapter{ТЕХНИЧЕСКАЯ ЧАСТЬ}

\section{Архитектура проекта}

Архитектура ChessView отражает фактическую структуру репозитория: `backend`, `frontend`, `admin-frontend`, `docs`, `docker-compose.yml`, `justfile`. Общая схема приведена на рисунке~\ref{fig:architecture}.

\begin{figure}[h]
\centering
\fbox{\begin{minipage}{0.86\textwidth}
\centering
Browser / React + Vite frontend\\
$\downarrow$ REST `/api/v1/*` \quad $\downarrow$ WebSocket `/ws?token=...`\\
FastAPI backend\\
$\downarrow$ SQLAlchemy async / Alembic\\
PostgreSQL\\[0.3cm]
Admin browser $\rightarrow$ admin-frontend $\rightarrow$ `/api/v1/admin/*`
\end{minipage}}
\caption{Общая архитектура ChessView}\label{fig:architecture}
\end{figure}

Backend организован по доменам: identity, game, matchmaking, profiles, ratings, communication, tournaments, scheduled\_matches, puzzles, payments, rtc и admin. Во многих доменах используется разделение на domain, application, infrastructure и presentation. Frontend организован по слоям app, pages, widgets, features, entities и shared.

\section{База данных и миграции}

Схема базы данных управляется Alembic. В репозитории присутствуют миграции `0001_baseline.py`--`0008_face_templates.py`. Логическая схема основных сущностей приведена на рисунке~\ref{fig:erd}.

\begin{figure}[h]
\centering
\fbox{\begin{minipage}{0.90\textwidth}
\centering
users $\rightarrow$ games $\rightarrow$ moves\\
games $\rightarrow$ chat\_messages\\
users $\rightarrow$ puzzle\_attempts $\leftarrow$ puzzles\\
tournaments $\rightarrow$ tournament\_players, tournament\_rounds, tournament\_pairings\\
scheduled\_matches $\rightarrow$ games\\
users $\rightarrow$ payment\_intents $\rightarrow$ payment\_events\\
users $\rightarrow$ face\_verification\_profiles, face\_verification\_sessions
\end{minipage}}
\caption{Логическая схема основных сущностей базы данных}\label{fig:erd}
\end{figure}

\section{REST API и WebSocket}

REST API регистрируется под префиксом `/api/v1`. Основные endpoints приведены в таблице~\ref{tab:api}.

\begin{center}
\scriptsize
\setlength{\tabcolsep}{4pt}
\renewcommand{\arraystretch}{1.25}
\begin{longtable}{|>{\RaggedRight\arraybackslash}p{0.14\linewidth}|>{\RaggedRight\arraybackslash}p{0.34\linewidth}|>{\RaggedRight\arraybackslash}p{0.34\linewidth}|>{\RaggedRight\arraybackslash}p{0.10\linewidth}|}
\caption{Основные REST endpoints}\label{tab:api}\\
\hline
Метод & Endpoint & Назначение & Доступ\\ \hline
\endfirsthead
\hline
Метод & Endpoint & Назначение & Доступ\\ \hline
\endhead
POST & `/api/v1/identity/register` & Регистрация пользователя & Public\\ \hline
POST & `/api/v1/identity/login` & Получение токенов & Public\\ \hline
GET & `/api/v1/identity/me` & Текущий пользователь & Auth\\ \hline
GET & `/api/v1/profiles/me` & Профиль текущего пользователя & Auth\\ \hline
GET & `/api/v1/games` & История партий & Auth\\ \hline
GET & `/api/v1/puzzles` & Список задач & Auth\\ \hline
POST & `/api/v1/tournaments` & Создание турнира & Auth\\ \hline
GET & `/api/v1/scheduled-matches/me` & Запланированные матчи & Auth\\ \hline
GET & `/api/v1/admin/users` & Список пользователей для администратора & Admin\\ \hline
\end{longtable}
\end{center}

WebSocket endpoint находится по адресу `/ws?token=<access_token>`. Основные события приведены в таблице~\ref{tab:ws}.

\begin{center}
\scriptsize
\setlength{\tabcolsep}{4pt}
\renewcommand{\arraystretch}{1.25}
\begin{longtable}{|>{\RaggedRight\arraybackslash}p{0.22\linewidth}|>{\RaggedRight\arraybackslash}p{0.22\linewidth}|>{\RaggedRight\arraybackslash}p{0.42\linewidth}|}
\caption{Основные WebSocket-события}\label{tab:ws}\\
\hline
Событие & Направление & Назначение\\ \hline
\endfirsthead
\hline
Событие & Направление & Назначение\\ \hline
\endhead
`queue\_join` & Клиент $\rightarrow$ сервер & Вход в очередь с выбранным контролем времени\\ \hline
`queue\_leave` & Клиент $\rightarrow$ сервер & Выход из очереди\\ \hline
`match\_found` & Сервер $\rightarrow$ клиент & Найдена пара и создана партия\\ \hline
`move` & Клиент $\rightarrow$ сервер & Отправка UCI-хода на проверку\\ \hline
`game\_state` & Сервер $\rightarrow$ клиент & Рассылка актуального состояния партии\\ \hline
`game\_over` & Сервер $\rightarrow$ клиент & Завершение партии\\ \hline
`chat\_send` / `chat\_message` & Оба направления & Чат партии\\ \hline
`rtc\_offer`, `rtc\_answer`, `rtc\_ice` & Оба направления & WebRTC signaling\\ \hline
\end{longtable}
\end{center}

\section{Игровой процесс}

Критическая логика live-партии находится на backend. Файл `backend/domains/game/domain/moves.py` создает шахматную доску из FEN, определяет ожидаемого игрока, разбирает UCI-ход, проверяет `board.legal_moves`, обновляет часы и сохраняет новый FEN. Такой подход снижает риск, что клиент изменит состояние партии произвольно.

Контроль времени реализован в доменной модели игры. При ходе создается clock snapshot, затем активному игроку начисляется increment. Фоновый монитор активных игр находится в `backend/domains/game/presentation/runtime.py`; он проверяет таймауты и disconnect grace. Эта реализация подходит для одного backend-процесса, но требует отдельной координации при горизонтальном масштабировании.

\section{Модули приложения}

Статус основных модулей приведен в таблице~\ref{tab:modules}.

\begin{center}
\scriptsize
\setlength{\tabcolsep}{4pt}
\renewcommand{\arraystretch}{1.25}
\begin{longtable}{|>{\RaggedRight\arraybackslash}p{0.18\linewidth}|>{\RaggedRight\arraybackslash}p{0.40\linewidth}|>{\RaggedRight\arraybackslash}p{0.28\linewidth}|}
\caption{Функциональные модули ChessView}\label{tab:modules}\\
\hline
Модуль & Фактическая реализация & Статус\\ \hline
\endfirsthead
\hline
Модуль & Фактическая реализация & Статус\\ \hline
\endhead
Авторизация & Identity router, JWT, password hashing, protected routes & Реализовано\\ \hline
Профили & Profile endpoints, leaderboard, search, head-to-head & Реализовано\\ \hline
Матчмейкинг & WebSocket queue, time controls, game creation & Реализовано, in-memory\\ \hline
Live game & Server-authoritative moves, clocks, reconnect, timeout, chat, RTC signaling & Реализовано\\ \hline
Анализ & Local Stockfish worker, FEN/PGN, board editor & Реализовано на frontend\\ \hline
Задачи & Puzzle catalog and attempts & Реализовано\\ \hline
Турниры & Lifecycle, standings, Swiss helpers, OTB type & Реализовано в базовом виде\\ \hline
Scheduled matches & Invite/accept/decline/cancel/start lifecycle & Реализовано\\ \hline
Payments & Payment intent emulator and coin debit/refund helpers & Эмулятор\\ \hline
Admin & Separate admin frontend and `/api/v1/admin/*` endpoints & Базовый admin scenario\\ \hline
Shop & Local inventory in browser, profile coins display & Демонстрационный UI\\ \hline
Clubs & Local-state club cards and create/search UI & Демонстрационный UI\\ \hline
Verification & Local face/passkey models, sessions, challenges & Архитектурная основа\\ \hline
\end{longtable}
\end{center}

\section{Проверка работоспособности}

Проверка проекта должна быть основана на фактически запускаемых командах. Сводная таблица приведена в таблице~\ref{tab:verification}. В графе фактического результата необходимо указывать результат последнего локального запуска.

\begin{center}
\scriptsize
\setlength{\tabcolsep}{3pt}
\renewcommand{\arraystretch}{1.20}
\begin{longtable}{|>{\RaggedRight\arraybackslash}p{0.19\linewidth}|>{\RaggedRight\arraybackslash}p{0.12\linewidth}|>{\RaggedRight\arraybackslash}p{0.28\linewidth}|>{\RaggedRight\arraybackslash}p{0.18\linewidth}|>{\RaggedRight\arraybackslash}p{0.13\linewidth}|}
\caption{План проверки проекта}\label{tab:verification}\\
\hline
Проверка & Тип & Команда или сценарий & Ожидаемый результат & Доказательство\\ \hline
\endfirsthead
\hline
Проверка & Тип & Команда или сценарий & Ожидаемый результат & Доказательство\\ \hline
\endhead
Backend compile & Авто & `python -m compileall app domains infrastructure shared tests` & Exit code 0 & terminal output\\ \hline
Backend tests & Авто & `pytest tests` & Tests passed & pytest summary\\ \hline
Migrations apply & Авто & `alembic upgrade head` & Current schema applied & Alembic output\\ \hline
Migration drift & Авто & `alembic check` & No new upgrade operations & Alembic output\\ \hline
Healthcheck & Smoke & `GET /health` & JSON status ok, version 1.0.1 & HTTP response\\ \hline
Protected API & Smoke & `GET /api/v1/puzzles` without token & HTTP 401 & HTTP status\\ \hline
Frontend lint & Авто & `yarn lint` & Exit code 0 & ESLint output\\ \hline
Frontend build & Авто & `yarn build` & TypeScript and Vite build complete & build output\\ \hline
Admin build & Авто & `yarn typecheck`, `yarn build` & Exit code 0 & terminal output\\ \hline
Docker config & Static & `docker compose config` & Valid Compose file & command output\\ \hline
Manual route check & Ручная & login/register, lobby, game, history, analysis & Pages open without critical errors & screenshots or demo\\ \hline
WebSocket check & Ручная & two sessions, `queue_join`, `move` & match and move events delivered & browser/devtools\\ \hline
\end{longtable}
\end{center}

В текущей версии корректно говорить о базовой функциональной проверке и регрессионных проверках. Нагрузочное тестирование, промышленная отказоустойчивость и проверка горизонтального масштабирования не заявляются, если они отдельно не проводились.

\section{Ограничения текущей реализации}

Ограничения не снижают ценность проекта, а показывают границы текущей инженерной версии. ChessView рассчитан на локальный или single-instance запуск. Matchmaking queue, WebSocket rooms и game monitor находятся в памяти backend-процесса. Uploaded media хранится локально. Для горизонтального масштабирования потребуются Redis/pub-sub, распределенная координация фонового монитора и object storage.

Платежный модуль является эмулятором и не интегрирован с реальным платежным провайдером. Shop и Clubs являются демонстрационными UI-модулями. Face/passkey verification содержит локальную архитектурную основу, но не является сертифицированной проверкой личности. WebSocket token в query string удобен для SPA-сценария, но требует дополнительной оценки безопасности для более строгих окружений.

\chapter{ОХРАНА ТРУДА}

\section{Требования безопасности при работе с вычислительной техникой}

Работа над программным проектом связана с длительным использованием персонального компьютера, монитора, клавиатуры, мыши и сетевого оборудования. Перед началом работы необходимо убедиться в исправности кабелей, блока питания, розеток и корпуса устройства. Запрещается работать с оборудованием, имеющим признаки перегрева, повреждения проводов или нестабильного питания.

При возникновении запаха гари, искрения, необычного шума или перегрева оборудование следует выключить безопасным способом и сообщить ответственному лицу. Самостоятельный ремонт электрической части без соответствующей квалификации недопустим.

\section{Организация рабочего места}

Рабочее место программиста должно обеспечивать устойчивую посадку, удобное положение рук и нормальный обзор монитора. Монитор рекомендуется располагать на расстоянии около 50--70 см от глаз, верхняя часть экрана должна находиться примерно на уровне глаз. Клавиатура и мышь должны находиться на одной высоте, чтобы снизить нагрузку на кисти.

Освещение должно быть равномерным, без сильных бликов на экране. При длительной работе необходимо делать короткие перерывы, менять позу, выполнять разминку для кистей, шеи и плеч, а также переводить взгляд на удаленные объекты.

\section{Вредные факторы и профилактика}

\begin{center}
\scriptsize
\setlength{\tabcolsep}{4pt}
\renewcommand{\arraystretch}{1.25}
\begin{longtable}{|>{\RaggedRight\arraybackslash}p{0.24\linewidth}|>{\RaggedRight\arraybackslash}p{0.34\linewidth}|>{\RaggedRight\arraybackslash}p{0.24\linewidth}|}
\caption{Вредные факторы при работе над проектом}\label{tab:safety}\\
\hline
Фактор & Причина & Профилактика\\ \hline
\endfirsthead
\hline
Фактор & Причина & Профилактика\\ \hline
\endhead
Зрительное напряжение & Длительная работа с кодом, браузером и PDF & Перерывы, настройка яркости, удобный масштаб\\ \hline
Статическая нагрузка & Долгое сидение в одной позе & Регулировка кресла, разминка, смена позы\\ \hline
Нагрузка на кисти & Набор текста и работа с мышью & Правильная высота клавиатуры, горячие клавиши\\ \hline
Психоэмоциональное напряжение & Сроки, отладка, параллельная работа с кодом и текстом & Планирование, чек-листы, короткие этапы проверки\\ \hline
Электробезопасность & Поврежденные кабели, перегрузка розеток & Проверка оборудования, аккуратная организация кабелей\\ \hline
\end{longtable}
\end{center}

\unnumberedchapter{ЗАКЛЮЧЕНИЕ}

В результате выполнения дипломного проекта разработано веб-приложение ChessView. Цель работы достигнута: создано full-stack приложение с React/TypeScript frontend, FastAPI backend, PostgreSQL, Alembic migrations, REST API, WebSocket endpoint и шахматной доменной логикой.

В проекте реализованы регистрация и вход, профили пользователей, история партий, матчмейкинг, серверная обработка ходов, контроль времени, reconnect/timeout logic, локальный анализ позиций, задачи, турниры, scheduled matches и административный интерфейс. Дополнительные разделы shop и clubs используются как демонстрационные расширенные сценарии, а платежи реализованы через эмулятор.

Практическая ценность ChessView заключается в том, что проект показывает полный цикл разработки учебной веб-платформы: архитектуру, UI, API, базу данных, realtime events, миграции, тесты и документацию запуска. Перспективами развития являются Redis-backed matchmaking and WebSocket fanout, object storage для медиа, расширенные E2E/load tests, развитие клубов, магазина, уведомлений, античит-механизмов и серверного анализа.

\unnumberedchapter{СПИСОК ИСПОЛЬЗОВАННЫХ ИСТОЧНИКОВ}

\begin{enumerate}[leftmargin=1.25cm,itemsep=0pt,topsep=0pt]
\item React. Documentation [Электронный ресурс]. — Режим доступа: \url{https://react.dev/} (дата обращения: 28.05.2026).
\item TypeScript. Documentation [Электронный ресурс]. — Режим доступа: \url{https://www.typescriptlang.org/docs/} (дата обращения: 28.05.2026).
\item Vite. Guide [Электронный ресурс]. — Режим доступа: \url{https://vite.dev/guide/} (дата обращения: 28.05.2026).
\item FastAPI. Documentation [Электронный ресурс]. — Режим доступа: \url{https://fastapi.tiangolo.com/} (дата обращения: 28.05.2026).
\item SQLAlchemy. Documentation [Электронный ресурс]. — Режим доступа: \url{https://docs.sqlalchemy.org/} (дата обращения: 28.05.2026).
\item Alembic. Documentation [Электронный ресурс]. — Режим доступа: \url{https://alembic.sqlalchemy.org/} (дата обращения: 28.05.2026).
\item PostgreSQL. Documentation [Электронный ресурс]. — Режим доступа: \url{https://www.postgresql.org/docs/} (дата обращения: 28.05.2026).
\item Uvicorn. Documentation [Электронный ресурс]. — Режим доступа: \url{https://www.uvicorn.org/} (дата обращения: 28.05.2026).
\item PyJWT. Documentation [Электронный ресурс]. — Режим доступа: \url{https://pyjwt.readthedocs.io/} (дата обращения: 28.05.2026).
\item python-chess. Documentation [Электронный ресурс]. — Режим доступа: \url{https://python-chess.readthedocs.io/} (дата обращения: 28.05.2026).
\item chess.js. Documentation [Электронный ресурс]. — Режим доступа: \url{https://www.npmjs.com/package/chess.js} (дата обращения: 28.05.2026).
\item react-chessboard. Documentation [Электронный ресурс]. — Режим доступа: \url{https://www.npmjs.com/package/react-chessboard} (дата обращения: 28.05.2026).
\item TanStack Query. Documentation [Электронный ресурс]. — Режим доступа: \url{https://tanstack.com/query/latest} (дата обращения: 28.05.2026).
\item Zustand. Documentation [Электронный ресурс]. — Режим доступа: \url{https://docs.pmnd.rs/zustand/getting-started/introduction} (дата обращения: 28.05.2026).
\item Docker Compose. Documentation [Электронный ресурс]. — Режим доступа: \url{https://docs.docker.com/compose/} (дата обращения: 28.05.2026).
\item FIDE Rules Commission. Laws of Chess [Электронный ресурс]. — Режим доступа: \url{https://rcc.fide.com/documentation/} (дата обращения: 28.05.2026).
\item Chess.com. Features [Электронный ресурс]. — Режим доступа: \url{https://www.chess.com/} (дата обращения: 28.05.2026).
\item Lichess. Free Online Chess [Электронный ресурс]. — Режим доступа: \url{https://lichess.org/} (дата обращения: 28.05.2026).
\item Tornelo. Online chess tournament platform [Электронный ресурс]. — Режим доступа: \url{https://home.tornelo.com/} (дата обращения: 28.05.2026).
\item MDN Web Docs. WebSocket API [Электронный ресурс]. — Режим доступа: \url{https://developer.mozilla.org/en-US/docs/Web/API/WebSocket} (дата обращения: 28.05.2026).
\item MDN Web Docs. WebRTC API [Электронный ресурс]. — Режим доступа: \url{https://developer.mozilla.org/en-US/docs/Web/API/WebRTC_API} (дата обращения: 28.05.2026).
\item OWASP Foundation. Authentication Cheat Sheet [Электронный ресурс]. — Режим доступа: \url{https://cheatsheetseries.owasp.org/cheatsheets/Authentication_Cheat_Sheet.html} (дата обращения: 28.05.2026).
\item International Labour Organization. Occupational safety and health [Электронный ресурс]. — Режим доступа: \url{https://www.ilo.org/topics/safety-and-health-work} (дата обращения: 28.05.2026).
\end{enumerate}

\unnumberedchapter{ПРИЛОЖЕНИЯ}

\appsection{Приложение А. Структура проекта}

\begin{lstlisting}
chessview/
  backend/
    app/
    domains/
    infrastructure/
    shared/
    alembic/
    tests/
  frontend/
    src/app
    src/pages
    src/widgets
    src/features
    src/entities
    src/shared
  admin-frontend/
  docs/
  docker-compose.yml
  justfile
  README.md
\end{lstlisting}

\appsection{Приложение Б. Фрагменты программного кода}

Фрагмент из `backend/domains/game/domain/moves.py`:

\begin{lstlisting}
board = chess.Board(game.fen)
expected_player = game.white_id if board.turn == chess.WHITE else game.black_id
if user_id != expected_player:
    raise NotYourTurn()

parsed_move = board.parse_uci(uci)
if parsed_move not in board.legal_moves:
    raise IllegalMove()

snapshot = capture_clock_snapshot(game, now)
board.push(parsed_move)
game.fen = board.fen()
apply_board_outcome(game, board, now)
\end{lstlisting}

Фрагмент из `backend/app/ws_entry.py`:

\begin{lstlisting}
data = json.loads(raw)
envelope = WSEnvelope(**data)
handler = EVENT_HANDLERS.get(envelope.type)
if handler is None:
    await manager.send_error(user_id, "INVALID_EVENT", ...)
    continue

async with async_session_factory() as session:
    await handler(envelope, user_id, session)
\end{lstlisting}

Фрагмент из `frontend/src/shared/api/http.ts`:

\begin{lstlisting}
if (token) {
  headers.Authorization = `Bearer ${token}`;
}

const fullUrl = buildUrl(API_BASE_URL, endpoint);
response = await fetch(fullUrl, { ...options, headers });
\end{lstlisting}

Фрагмент из `frontend/src/shared/api/ws.ts`:

\begin{lstlisting}
const envelope: WSEnvelope<EventPayload<K>> = {
  type,
  payload,
  game_id: gameId,
  timestamp: new Date().toISOString(),
};
this.ws.send(JSON.stringify(envelope));
\end{lstlisting}

\appsection{Приложение В. Примеры API-запросов}

\begin{lstlisting}
POST /api/v1/identity/register
Content-Type: application/json

{
  "username": "player",
  "email": "player@example.com",
  "password": "Password123!"
}
\end{lstlisting}

\begin{lstlisting}
GET /api/v1/games?page=1&size=20
Authorization: Bearer <access_token>
\end{lstlisting}

\begin{lstlisting}
Endpoint: /ws?token=<access_token>
{
  "type": "move",
  "game_id": "<uuid>",
  "payload": { "uci": "e2e4" },
  "timestamp": "2026-05-28T00:00:00Z"
}
\end{lstlisting}

\appsection{Приложение Г. Материалы для финальной демонстрации}

Перед сдачей рекомендуется обновить актуальные скриншоты из локально запущенного приложения: главная страница, форма входа, dashboard, лобби, live-партия, анализ, задачи, турниры, профиль, scheduled matches, OTB manager и admin frontend. Скриншоты должны соответствовать последнему коду и не должны заменяться макетами.

\clearpage
\thispagestyle{plain}
\begin{center}
ОТЗЫВ

на дипломный проект
\end{center}

Студента: \filledfield{9cm}{Сакенов Ануар Сункарулы}

Группы: \filledfield{9cm}{ПО 2302}

Тема дипломного проекта: «Разработка веб-приложения ChessView».

Руководитель: \filledfield{8cm}{Мартынцов Николай Викторович}

\vspace{0.5cm}
Текст отзыва заполняется руководителем.

\vspace{0.4cm}
Оценка руководителя: \blankfield{5cm}

\vspace{0.25cm}
\begin{flushright}
Руководитель: \blankfield{6cm}\\[0.4cm]
\filledfield{6cm}{Мартынцов Николай Викторович}\\[0.4cm]
Дата: \blankfield{7cm}
\end{flushright}

\clearpage
\thispagestyle{plain}
\begin{center}
РЕЦЕНЗИЯ

на дипломный проект
\end{center}

Студента: \filledfield{9cm}{Сакенов Ануар Сункарулы}

Группы: \filledfield{9cm}{ПО 2302}

Тема дипломного проекта: «Разработка веб-приложения ChessView».

Рецензент: \blankfield{9cm}

\vspace{0.5cm}
Текст рецензии заполняется рецензентом.

\vspace{0.4cm}
Оценка рецензента: \blankfield{5cm}

\vspace{0.25cm}
\begin{flushright}
Рецензент: \blankfield{6cm}\\[0.4cm]
\blankfield{6cm}\\[0.4cm]
Дата: \blankfield{7cm}
\end{flushright}

\end{document}
